\scanpage{32}
INSTITUT\\
DES HAUTES ÉTUDES SCIENTIFIQUES\\
BURES-SUR-YVETTE (S.-\&-O.)\\
928-48-53

\uline{Corps de base topologique isom. \unsure{à} $\mathbf R$ ou $\mathbf C$}

$K$ corps topologique isom. \unsure{à} $\mathbf C$ (il y a donc \uline{deux} isom avec $\mathbf C$, mais on ne choisit pas). $K$ contient can. $\mathbf R$.

$M$ motif sur $K$. Alors \deleted{$\operatorname{DR}(M)$ \unsure{De Rham}} pour \unsure{tt} $X$ lisse projectif sur $K$, $X(K)$ est un espace topologique cpt (variété $K$-analytique compacte), d'où $H^*(X,\mathbf Q)$, et on définit ainsi \uline{foncteur de Betti} $M\rightsquigarrow\operatorname{Betti}(M)$ sur les motifs sur $K$. $\operatorname{Betti}(M)\otimes_{\mathbf Q}K$ $\unsure{\simeq}$ canoniquement à $H^*(X,K)$. Par l'argument habituel, ce dernier est isomorphe can. à la coh. de De R., d'où
\[
\operatorname{Betti}(M)_K\simeq\operatorname{DR}(M)
\qquad\otimes\text{-isom can.}
\]
Ceci met sur \inserted{(ce $K$-vecteur)} $\operatorname{DR}(M)$ une structure \uline{réelle}, dont le sous-espace réel est $\operatorname{Betti}(M)_{\mathbf R}$ ; soit $\sigma$ l'involution. La filtration canonique de $\operatorname{DR}(M)$, et $\sigma$, permettent de façon unique de splitter cette filtration, de façon à obtenir sur $\operatorname{DR}(M)$ une \uline{bigraduation} compatible avec les degrés \unsure{totaux}, et \unsure{redonnant} la filtration can. à l'aide du premier degré. Le choix d'un isom $\unsure{\pi|i}:K\unsure{-|\longrightarrow}\mathbf C$ met, par transport de str., une

\mspage{31}
